\section{AgentGS}
\label{sec:method}
AgentGS organizes scientific tool use around reusable operations and their dependencies. Given a scientific question and available resources, the agent retrieves relevant strategies, composes their local operations into a task-specific plan, completes missing connections, and grounds the result in concrete tool calls. Figure~\ref{fig:overview} summarizes this workflow. Its central design choice is to preserve the strategy-specific conditions under which an operation can be reused while allowing that operation to participate in a new composition.

\subsection{Task Formulation and Design}
\label{sec:task-formulation}
Let $q$ be a natural-language task, $X$ its available input resources, and $\mathcal{T}$ the registered tool collection. A task can request both computed quantities and output artifacts: for example, a composite raster, a statistic calculated from that raster, and a comparison with a supplied reference. Solving the task requires determining the necessary operations, their data dependencies, admissible implementations, and concrete arguments.

A useful plan must meet three conditions. First, it must cover the requested operations and their prerequisites. Second, each operation must have an available implementation whose inputs can be supplied by $X$ or by successful preceding operations. Third, the produced quantities and artifacts must support the requested answer. The first two conditions constrain workflow construction; the third requires checking the meaning of the actual results. A structurally valid plan alone does not establish scientific correctness.

AgentGS separates this problem into capability planning and tool execution. Capability planning determines what operations are needed and how they connect. Tool grounding chooses implementations under the resulting constraints, and execution resolves their arguments from observed data. This allows alternative implementations of an operation without discarding its role in the workflow. It also makes a missing intermediate operation distinguishable from a failed invocation of an already selected operation.

\subsection{Skill--Capability--Tool Representation}
\label{sec:representation}
The representation distinguishes strategies, shared operations, and executable interfaces (Figure~\ref{fig:overview}(b)).

\paragraph{Skills as reusable strategies.}
A Skill describes a class of workflows through an objective, applicable inputs, expected artifacts, capability steps, and usage constraints. Each step specifies its operation, mandatory predecessors, required or optional status, and permitted implementations. Validation requirements and recovery policies describe how the strategy's results and failures should be handled. Reuse selects task-relevant local steps; an agent need not execute every operation in the original strategy.

\paragraph{Capabilities as shared operations.}
A Capability specifies an operation through its function, input/output types, preconditions, postconditions, and candidate tools. The same capability can occur in several Skills. Its occurrence retains the source strategy's constraints: two Skills may both use raster analysis but permit different tools or require different preceding transformations. Sharing a capability therefore provides a point of connection between strategies without making all its occurrences interchangeable.

\paragraph{Tools as implementations.}
A Tool is a callable interface with an identifier, parameter schema, input/output definitions, and runtime information. Capability--Tool mappings record which tools implement an operation and any operation-specific parameter constraints. The mapping is many-to-many: a capability can admit several implementations, and a tool can support several capabilities. Concrete files, values, and upstream artifacts are bound at execution time.

\paragraph{Skill hypergraph.}
Let $S$ and $C$ denote the registered Skills and Capabilities. Their membership defines
\begin{equation}
H=(C,E),\qquad E=\{e_s:s\in S\},\qquad e_s\subseteq C.
\label{eq:hypergraph}
\end{equation}
Capabilities are nodes and each Skill is a hyperedge. Skills containing $\{a,b\}$ and $\{b,c\}$ intersect at capability $b$, providing a candidate connection for composition. Hyperedge membership records which operations belong to a strategy. It does not, by itself, specify the order or data flow between them.

Execution prerequisites are represented separately by a scoped relation $d=(s,U,c)$: within Skill $s$, capability $c$ requires every predecessor in $U$. This is a conjunctive requirement. When an operation requires both a prepared raster and a reference value, supplying either input alone does not satisfy the declaration. Conditional order is used only when its stated condition applies. The strategy identity remains part of each relation, preventing a dependency from one Skill from being silently substituted for another's.

\paragraph{Capability instances and artifact bindings.}
A task-specific plan is
\begin{equation}
P=(V,D_d,D_c,\beta),\qquad v_i=(c_i,s_i,r_i),
\label{eq:plan}
\end{equation}
where $c_i$ is the operation, $s_i$ its source Skill, and $r_i$ its step/input role. $D_d$ and $D_c$ represent data and control dependencies, respectively; $\beta$ binds inputs to task resources or outputs of preceding instances. A capability is a reusable definition, whereas an instance is one occurrence on particular inputs. Two operations on different rasters remain distinct even if they select the same tool. This distinction is essential for preserving multiplicity, provenance, and the identity of the artifact consumed downstream.

\subsection{Plan Composition}
\label{sec:composition}
\label{sec:initial-selection}
Composition turns retrieved strategies into a task-specific set of operations. Skill retrieval uses the question, domain, and input types to supply candidate strategies. Selection then identifies applicable Skills and the operations needed for the requested outputs. The LLM makes these decisions within the retrieved declarations, retaining the relevant source Skills and their constraints.

\paragraph{Mandatory-predecessor expansion.}
Selected operations can depend on outputs not yet included in the plan. For a selected Skill $s$, let $R_s^{(0)}$ contain its initially required capabilities. Its dependency closure can be written as
\begin{equation}
R_s^{(k+1)}
=
R_s^{(k)}
\cup
\bigcup_{\substack{(s,U,c)\in\mathcal{D}\\c\in R_s^{(k)}}} U,
\label{eq:closure}
\end{equation}
where $\mathcal{D}$ is the set of declared mandatory dependencies. Expansion continues until no required predecessor is missing. This expression describes closure within a strategy; instantiation still preserves separate occurrences and input roles. Missing declarations and dependency cycles are diagnosed rather than accepted as executable plans.

\paragraph{Combining local fragments.}
Composition assembles the selected operations and their required predecessors into $P$. Composition retains the provenance and admissible tool scope of every occurrence, together with its data and control dependencies. Shared capability identity makes reuse discoverable, but does not justify merging operations on different inputs or changing a strategy's prerequisites.

Predecessor expansion only supplies dependencies already declared for the selected operations. It can leave a gap between fragments from different strategies. For instance, a compositing fragment may produce a raster while a comparison fragment expects a scalar statistic. Both fragments can be internally complete and still fail to form an executable analysis. Completion addresses these remaining gaps.

When no Skill is applicable, AgentGS retrieves tool candidates from GeoToolKG for direct tool-call planning. Both routes use the same execution and answer-generation interfaces. This fallback provides access to registered tools, but it does not supply the strategy knowledge absent from retrieval.

\subsection{Automatic Plan Completion}
\label{sec:completion}
AgentGS examines the initial plan for omitted operations and disconnected fragments. It proposes supporting operations from compatible Skills for task-relevance review, then searches other Skills for declared branches connecting active fragments. The search fixes the active boundaries and uses shared operations and typed interfaces to identify candidate connections within individual provider Skills.

A candidate must add a missing operation, satisfy its prerequisites, and have available tools within the declared scope. Insertion must preserve acyclic dependencies. A unique valid branch is added automatically; ambiguous alternatives are reported. This procedure is designed to repair omissions without requiring the model to generate every intermediate step.

Insertion preserves the entry operation's constraints and updates the receiving operation's dependencies and bindings. Recomposition verifies mandatory prerequisites; a branch requiring additional operations outside the accepted connection is rejected.

Figure~\ref{fig:pipeline} illustrates a connection between raster compositing and scalar comparison. Completion inserts a valid-pixel mean operation and binds the comparison to its output. Grounding selects concrete tools and preserves these producer--consumer bindings; an inconsistent grounding is rejected. The resulting plan uses the existing execution and answer-generation services. Search rules, execution details, and the complete example are provided in Appendices~\ref{app:method-details}, \ref{sec:system}, and~\ref{sec:case-study}.

\subsection{Grounding, Execution, and Verification}
\label{sec:execution}
A completed capability plan describes operations and dependencies. Grounding converts it into execution instances with concrete tool identifiers, arguments, and references to upstream results.

\paragraph{Constrained tool selection.}
For each capability instance, the candidate set is restricted by its source Skill, operation requirements, tool availability, and interface constraints. GeoToolKG supplies tool descriptions and associations for retrieval; those associations do not override the plan's dependencies. The LLM selects implementations within the admitted scope. Required operations and scoped ordering are checked again after the concrete chain is produced.

\paragraph{Explicit input propagation.}
Each execution instance has an identity and declares its predecessors. Input bindings name a producing instance and a field in that instance's observed output. At execution time, the system resolves these references and validates the resulting arguments against the selected tool's schema. A consumer cannot substitute an unrelated earlier call merely because it used the same tool name. Unknown predecessors, unresolved references, and incompatible arguments prevent invocation.

This separation is especially important for scientific artifacts. In a compositing task, the statistics call must consume the generated composite, not another raster mentioned in the question. Likewise, a comparison must use the statistic actually returned by its producer. The execution record links these inputs to their sources, making artifact use inspectable independently of the final answer.

\paragraph{Failure handling.}
Execution failures produce structured feedback for bounded argument repair or replanning when recovery is enabled. Literal arguments may be corrected, while observed upstream bindings remain authoritative. Replanning validates instance identities, dependencies, and scopes again. Reuse of a successful result requires unchanged inputs and instance semantics; it is not inferred from a common tool-name prefix. The runtime conservatively reruns file-backed and external-data computations across plan attempts, while allowing reuse of explicitly supported pure scalar computations when their complete dependencies are verified.

Tool recovery, initial plan repair, and transport retries have separate roles and budgets. Disabling execution recovery terminates the scientific chain after an execution failure while retaining initial planning and input validation. Engineering details and the distinction between these budgets appear in Appendix~\ref{sec:system}.

\paragraph{Verification and answer generation.}
The answer is constructed from observed results and artifacts. Verification checks whether the requested outputs are supported and records unfinished operations when execution is incomplete. Successful invocation, evidence consistency, and agreement with a benchmark reference are distinct properties; the benchmark evaluates the last of these independently.

A verified composition can also be represented as a candidate reusable Skill, preserving its input roles, dependencies, and usage constraints. This is a strategy representation over existing operations, not an expansion of the underlying tool collection. The present results do not establish a performance or cost benefit from autonomous Skill accumulation; the candidate-registration design is described in Appendix~\ref{app:skill-reuse}.

